\documentclass[a4paper,11pt]{article}
\usepackage[utf8]{inputenc}
\usepackage[T1]{fontenc}
\usepackage[english]{babel}
\usepackage[left=1cm,right=1cm,top=.5cm,bottom=0cm]{geometry}
\usepackage{enumitem}
\usepackage{hyperref}
\usepackage{xcolor}
\usepackage{titlesec}
\usepackage{marvosym}
\usepackage{lmodern}
\usepackage[normalem]{ulem}

\definecolor{primary}{RGB}{0, 51, 102}
\definecolor{secondary}{RGB}{80, 80, 80}
\hypersetup{colorlinks=true, linkcolor=primary, filecolor=primary, urlcolor=primary}
\titleformat{\section}{\large\bfseries\color{primary}\uppercase}{}{0em}{}[\titlerule]
\titlespacing{\section}{0pt}{8pt}{5pt}
\newcommand{\resumeItem}[1]{\item\small{#1}}

\begin{document}

\begin{center}
    {\Large \textbf{Lo\"ic Aron MBASSI EWOLO}} \\ \vspace{4pt}
    \textbf{\large Data Scientist -- Bureau d'Etudes} \\
    \textit{\small Stage INRIA Saclay (Avril--Aout 2026) -- Disponible en alternance des septembre 2026 (12 mois)} \\ \vspace{4pt}
    \small
    \Pointinghand\ Cergy, France (Mobile IDF) \quad
    \Mobilefone\ (+33) 7 59 52 33 98 \quad
    {\color{primary}\Letter\ \href{mailto:loicaron.mbassiewolo@ensea.fr}{loicaron.mbassiewolo@ensea.fr}} \\ \vspace{2pt}
    {\color{primary}LinkedIn: \href{https://www.linkedin.com/in/loic-mbassi/}{linkedin.com/in/loic-mbassi}} \quad
    {\color{primary}GitHub: \href{https://github.com/Nameless0l}{github.com/Nameless0l}} \quad
    {\color{primary}Portfolio: \href{https://mbassiloic.tech}{mbassiloic.tech}}
\end{center}

\vspace{-6pt}

\noindent\small{Data scientist with background in signal processing and deep learning for complex physical systems. Experienced in designing predictive models from sensor data, handling time-series analysis, implementing ML solutions for engineering applications, and bridging theoretical models with practical implementations.}

\section{Formation}
\begin{itemize}[leftmargin=*, label=\tiny$\bullet$, nosep]
    \item \textbf{ENSEA -- Ecole Nationale Superieure de l'Electronique} \hfill \textit{2025 -- 2027} \\
    \small{Double diplome d'Ingenieur -- Systemes Embarques, Signal, IA.}
    \item \textbf{Ecole Polytechnique de Yaounde} \hfill \textit{2021 -- 2025} \\
    \small{Diplome d'Ingenieur de Conception (Master 2) : Mathematiques Appliquees, Genie Logiciel.}
\end{itemize}

\section{Competences Techniques}
\begin{itemize}[leftmargin=*, nosep]
    \item \small{\textbf{Data Science \& ML:} TensorFlow, scikit-learn, pandas, numpy, deep learning, time-series}
    \item \small{\textbf{Signal Processing:} MATLAB/Simulink, filtrage, feature extraction, fourier analysis}
    \item \small{\textbf{Modeling \& Simulation:} MATLAB/Simulink, system dynamics, control systems, FDI}
    \item \small{\textbf{Data Pipeline:} Python, SQL, pandas, data preprocessing, feature engineering}
    \item \small{\textbf{Tools:} Git, Docker, Jupyter, scikit-learn, visualization (matplotlib, seaborn)}
\end{itemize}

\section{Experiences / Projets}
\begin{itemize}[leftmargin=*, label={-}]

\item \textbf{Stage Recherche -- INRIA Saclay, Equipe TRIBE} \hfill \textit{Avril -- Aout 2026} \\
\textit{Pipeline NLP et data science production} \hfill \textit{Python, pandas, BERT, scikit-learn}
\vspace{-2pt}
    \begin{itemize}[leftmargin=0.4cm, nosep, label=\tiny$\bullet$]
        \item \small{Conception pipeline NLP analyse a grande echelle avec extraction features et preprocessing}
        \item \small{Evaluation modeles statistiques pour scoring confiance et analyse risques}
    \end{itemize}
\vspace{3pt}

\item \textbf{Projet UROP 2024 -- Analyse Electrocardiogrammes} \hfill \textit{2024} \\
\textit{Deep learning pour imagerie medicale} \hfill \textit{TensorFlow/Keras, signal, traitement image}
\vspace{-2pt}
    \begin{itemize}[leftmargin=0.4cm, nosep, label=\tiny$\bullet$]
        \item \small{Modele detection et classification anomalies cardiaques a partir images ECG}
        \item \small{Extraction features signaux temporels et collaboration mentors Google DeepMind}
    \end{itemize}
\vspace{3pt}

\item \textbf{Fault Detection et Fault-Tolerant Control -- Drone} \hfill \textit{2024} \\
\textit{Systemes controle et detection pannes} \hfill \textit{MATLAB/Simulink, observateurs, FDI}
\vspace{-2pt}
    \begin{itemize}[leftmargin=0.4cm, nosep, label=\tiny$\bullet$]
        \item \small{Modelisation dynamique drone et conception detecteur pannes basee observateurs}
        \item \small{Strategie controle tolerant aux defauts avec reconfiguration controleur et simulations realistes}
    \end{itemize}
\vspace{3pt}

\item \textbf{Urban Waste Management -- Prediction et Optimisation} \hfill \textit{2024} \\
\textit{ML pour smart city} \hfill \textit{ML, IoT data, algorithmes optimisation}
\vspace{-2pt}
    \begin{itemize}[leftmargin=0.4cm, nosep, label=\tiny$\bullet$]
        \item \small{Modele prediction besoins collecte et algorithme optimisation routes (reduction 20\% couts)}
        \item \small{1er prix competition CITS contre 75 equipes}
    \end{itemize}

\end{itemize}

\end{document}